%% ============================================================
%%  PREAMBLE — Production-Grade Agentic AI: LangGraph + Python
%%  Extracted from main.tex.  Input ONCE from Tex/main.tex.
%% ============================================================

\documentclass[
  11pt,
  twoside,
  openright
]{book}

\usepackage{polyglossia}
\setdefaultlanguage{english}

\usepackage{fontspec}
\setmainfont{TeX Gyre Termes}[
  BoldFont   = TeX Gyre Termes Bold,
  ItalicFont = TeX Gyre Termes Italic
]
\setsansfont{TeX Gyre Heros}[
  BoldFont   = TeX Gyre Heros Bold,
  ItalicFont = TeX Gyre Heros Italic
]
\setmonofont{TeX Gyre Cursor}[
  Scale    = 0.88,
  BoldFont = TeX Gyre Cursor Bold
]

\usepackage{geometry}
\geometry{
  papersize  = {7in, 10in},
  top        = 1in,
  bottom     = 1.1in,
  inner      = 1in,
  outer      = 0.875in,
  headheight = 23.07pt,
  twoside
}

%% Fix fancyhdr \headheight warning (must be after geometry)
\setlength{\headheight}{23.07004pt}
\addtolength{\topmargin}{-0.00003pt}

\usepackage{xcolor}
\definecolor{BookBlue}       {HTML}{1A4A8A}
\definecolor{BookDark}       {HTML}{1C1C1E}
\definecolor{BookMid}        {HTML}{4A4A4A}
\definecolor{LearnBg}        {HTML}{EEF7EE}
\definecolor{LearnBorder}    {HTML}{27AE60}
\definecolor{TipBg}          {HTML}{EAF4FB}
\definecolor{TipBorder}      {HTML}{2980B9}
\definecolor{WarningBg}      {HTML}{FEF9E7}
\definecolor{WarningBorder}  {HTML}{F39C12}
\definecolor{NoteBg}         {HTML}{F4F4F4}
\definecolor{NoteBorder}     {HTML}{888888}
\definecolor{TakeawayBg}     {HTML}{F0F0FF}
\definecolor{TakeawayBorder} {HTML}{5B4FBE}
\definecolor{CodeBg}         {HTML}{F7F7F7}
\definecolor{CodeComment}    {HTML}{6A737D}
\definecolor{CodeKeyword}    {HTML}{D73A49}
\definecolor{CodeString}     {HTML}{032F62}
\definecolor{CodeConstant}   {HTML}{005CC5}
\definecolor{CodeAnnotation} {HTML}{6A737D}

\usepackage{microtype}
\usepackage{setspace}
\setstretch{1.18}

\usepackage{csquotes}
\usepackage{parskip}
\setlength{\parskip}{6pt}
\setlength{\parindent}{0pt}

\widowpenalty=10000
\clubpenalty=10000
\displaywidowpenalty=10000
\predisplaypenalty=10000
\raggedbottom
\emergencystretch=3em
\hbadness=3000
\vbadness=3000

\usepackage{needspace}
\usepackage{etoolbox}
\preto\section{\needspace{5\baselineskip}}
\preto\subsection{\needspace{4\baselineskip}}
\preto\subsubsection{\needspace{3\baselineskip}}

\usepackage{titlesec}
\renewcommand{\bottomtitlespace}{0.15\textheight}

\titleformat{\part}[display]
  {\centering\sffamily\Huge\bfseries\color{BookBlue}}
  {\sffamily\large\mdseries\partname\ \thepart}{20pt}{}

\titleformat{\chapter}[display]
  {\sffamily\huge\bfseries\color{BookBlue}}
  {\sffamily\normalsize\mdseries\chaptertitlename\ \thechapter}{12pt}{}
\titlespacing*{\chapter}{0pt}{0pt}{24pt}

\titleformat{\section}
  {\sffamily\Large\bfseries\color{BookBlue}}
  {\thesection}{1em}{}
\titlespacing*{\section}{0pt}{18pt}{6pt}

\titleformat{\subsection}
  {\sffamily\large\bfseries\color{BookDark}}
  {\thesubsection}{1em}{}
\titlespacing*{\subsection}{0pt}{14pt}{4pt}

\titleformat{\subsubsection}
  {\sffamily\normalsize\bfseries\color{BookDark}}
  {\thesubsubsection}{1em}{}
\titlespacing*{\subsubsection}{0pt}{10pt}{2pt}

\usepackage{fancyhdr}
\pagestyle{fancy}
\fancyhf{}
\fancyhead[LE]{\small\sffamily\color{BookMid}\leftmark}
\fancyhead[RO]{\small\sffamily\color{BookMid}\rightmark}
\fancyfoot[LE,RO]{\small\sffamily\color{BookMid}\thepage}
\renewcommand{\headrulewidth}{0.3pt}
\renewcommand{\footrulewidth}{0pt}

\fancypagestyle{plain}{
  \fancyhf{}
  \fancyfoot[LE,RO]{\small\sffamily\color{BookMid}\thepage}
  \renewcommand{\headrulewidth}{0pt}
}

\makeatletter
\renewcommand{\cleardoublepage}{%
  \clearpage
  \if@twoside
    \ifodd\c@page\else
      \thispagestyle{empty}\hbox{}\newpage
      \if@twocolumn\hbox{}\newpage\fi
    \fi
  \fi
}
\makeatother

\usepackage{rotating}
\usepackage{booktabs}
\usepackage{tabularx}
\usepackage{array}
\usepackage{longtable}
\renewcommand{\arraystretch}{1.35}

\newcolumntype{P}[1]{>{\RaggedRight\arraybackslash}p{#1}}

\usepackage{graphicx}
\usepackage{float}
\usepackage{caption}
\captionsetup{
  font      = {small,sf},
  labelfont = {bf,color=BookBlue},
  skip      = 6pt
}

\usepackage{listings}
\lstset{
  backgroundcolor   = \color{CodeBg},
  basicstyle        = \ttfamily\fontsize{9}{11}\selectfont,
  keywordstyle      = \color{CodeKeyword}\bfseries,
  stringstyle       = \color{CodeString}\itshape,
  commentstyle      = \color{CodeComment}\itshape,
  numbers           = left,
  numberstyle       = \tiny\color{BookMid},
  numbersep         = 6pt,
  breaklines        = true,
  frame             = single,
  framerule         = 0.4pt,
  rulecolor         = \color{NoteBorder},
  showstringspaces  = false,
  tabsize           = 2,
  captionpos        = b,
  breakatwhitespace = false,
  columns           = fullflexible,
  linewidth         = \linewidth,
  postbreak         = \mbox{\textcolor{gray}{$\hookrightarrow$}\space},
  xleftmargin       = 1.5em,
  framexleftmargin  = 1.5em,
  aboveskip         = 6pt,
  belowskip         = 4pt,
}

\lstdefinelanguage{Python}{
  morekeywords={
    False,None,True,and,as,assert,async,await,break,class,continue,def,
    del,elif,else,except,finally,for,from,global,if,import,in,is,lambda,
    nonlocal,not,or,pass,raise,return,try,while,with,yield
  },
  morekeywords=[2]{self,cls},
  sensitive=true,
  morecomment=[l]{\#},
  morestring=[b]",
  morestring=[b]',
  morestring=[s]{"""}{"""}, 
  morestring=[s]{'''}{'''},
  keywordstyle=\color{CodeKeyword}\bfseries,
  keywordstyle=[2]\color{CodeConstant},
  stringstyle=\color{CodeString},
  commentstyle=\color{CodeComment}\itshape,
  breaklines=true,
  breakatwhitespace=false,
  columns=fullflexible,
  keepspaces=true,
  showstringspaces=false,
  postbreak=\mbox{\textcolor{gray}{$\hookrightarrow$}\space},
}

\lstdefinelanguage{yaml}{
  morecomment=[l]{\#},
  morestring=[b]",
  morestring=[b]',
  literate=
    *{true}{{\color{CodeKeyword}true}}{4}
    {false}{{\color{CodeKeyword}false}}{5}
    {null}{{\color{CodeKeyword}null}}{4},
  sensitive=true,
  keywordstyle=\color{CodeKeyword}\bfseries,
  stringstyle=\color{CodeString},
  commentstyle=\color{CodeComment}\itshape,
  breaklines=true,
  columns=fullflexible,
  showstringspaces=false
}

\lstdefinelanguage{json}{
  morestring=[b]",
  morecomment=[l]//,
  morecomment=[s]{/*}{*/},
  sensitive=true,
  stringstyle=\color{CodeString},
  commentstyle=\color{CodeComment}\itshape,
}

\DeclareRobustCommand{\code}[1]{\texttt{#1}}
\providecommand{\term}[1]{\textit{#1}}
\providecommand{\chapterquote}[2]{%
  \vspace{1em}
  \begin{quote}\itshape#1\end{quote}
  \begin{flushright}--- #2\end{flushright}
  \vspace{1em}
}

\newcommand{\githubref}[1]{%
  \vspace{2pt}%
  {\small\sffamily\color{BookMid}Full source: \texttt{#1}}%
  \vspace{4pt}%
}

\usepackage{tcolorbox}
\tcbuselibrary{skins,breakable}

\newtcolorbox{learningbox}{
  colback=LearnBg, colframe=LearnBorder,
  before skip=12pt, fonttitle=\sffamily\bfseries\small,
  title={Learning Objective}, breakable, lines before break=3,
  left=8pt, right=8pt, top=6pt, bottom=6pt,
  boxrule=0.6pt, width=\linewidth
}
\newtcolorbox{tipbox}{
  colback=TipBg, colframe=TipBorder,
  before skip=12pt, fonttitle=\sffamily\bfseries\small,
  title={Tip}, breakable, lines before break=3,
  left=8pt, right=8pt, top=6pt, bottom=6pt,
  boxrule=0.6pt, width=\linewidth
}
\newtcolorbox{warningbox}{
  colback=WarningBg, colframe=WarningBorder,
  before skip=12pt, fonttitle=\sffamily\bfseries\small,
  title={Warning}, breakable, lines before break=3,
  left=8pt, right=8pt, top=6pt, bottom=6pt,
  boxrule=0.6pt, width=\linewidth
}
\newtcolorbox{notebox}{
  colback=NoteBg, colframe=NoteBorder,
  before skip=12pt, fonttitle=\sffamily\bfseries\small,
  title={Note}, breakable, lines before break=3,
  left=8pt, right=8pt, top=6pt, bottom=6pt,
  boxrule=0.6pt, width=\linewidth
}
\newtcolorbox{takeawaybox}[1][Key Takeaway]{
  colback=TakeawayBg, colframe=TakeawayBorder,
  before skip=12pt, fonttitle=\sffamily\bfseries\small,
  title={#1}, breakable, lines before break=3,
  left=8pt, right=8pt, top=6pt, bottom=6pt,
  boxrule=0.6pt, width=\linewidth
}

%% ── Source-Layout tree environment ─────────────────────────────────────
%% Usage:
%%   \begin{sourcelayout}
%%     \dirtree{%
%%       .1 src/.
%%       .2 agent/.
%%       .3 nodes/.
%%       .4 perceive\_node.py \DTcomment{§1.2}.
%%     }
%%   \end{sourcelayout}
%%
%% Requires: \usepackage{dirtree}  (loaded below)
%% Colours reuse NoteBg / NoteBorder (neutral grey — reference material,
%% not instructional content).
%% ────────────────────────────────────────────────────────────────────────
\usepackage{dirtree}

\newtcolorbox{sourcelayout}{%
  colback      = NoteBg,
  colframe     = NoteBorder,
  colbacktitle = NoteBorder,
  coltitle     = white,
  fonttitle    = \sffamily\bfseries\small,
  title        = {Source Layout},
  before skip  = 12pt,
  after skip   = 10pt,
  breakable,
  lines before break = 3,
  left   = 6pt,
  right  = 6pt,
  top    = 6pt,
  bottom = 6pt,
  boxrule = 0.6pt,
  width   = \linewidth,
}

%% dirtree global style: monospaced font, consistent with code listings
\renewcommand*\DTstylecomment{\rmfamily\color{CodeComment}\itshape}
\renewcommand*\DTstyle{\ttfamily\color{BookDark}}

\usepackage{amsmath}
\usepackage{amssymb}

\PassOptionsToPackage{hyphens}{url}
\usepackage{hyperref}
\usepackage{imakeidx}
\makeindex[columns=2, title=Index, intoc=true]
\urlstyle{tt}
\makeatletter
\g@addto@macro\UrlBreaks{\do\.\do\_\do\-}
\makeatother

\hypersetup{
  colorlinks  = true,
  linkcolor   = BookBlue,
  urlcolor    = BookBlue,
  citecolor   = BookBlue,
  pdftitle    = {Production-Grade Agentic AI: LangGraph + Python Frameworks},
  pdfauthor   = {Haythem Aber},
  pdfsubject  = {LangGraph, Python, Agentic AI, Production Systems},
  pdfkeywords = {LangGraph, Python, AI agents, production, observability, security}
}

\usepackage{tikz}
\usetikzlibrary{arrows.meta, positioning, shapes.geometric, shapes.misc, calc}
\usepackage{pdflscape}
\usepackage{ragged2e}

\newcommand{\component}[1]{\textsf{\textbf{#1}}}
\newcommand{\seechapter}[1]{\textit{(See Chapter~\ref{#1})}}
\newcommand{\seepart}[1]{\textit{(See Part~\ref{#1})}}

\newcommand{\partclose}[1]{%
  \cleardoublepage
  \vspace*{\fill}
  \begin{center}\sffamily\Large\color{BookBlue}#1\end{center}
  \vspace*{\fill}
  \cleardoublepage
}

%% Shared book-wide infrastructure layered on top of the author's preamble.
%% Keep this file loaded once by Tex/main.tex; chapter fragments must not load it.

\usepackage{calc} % Pandoc-generated longtable column dimensions.
\usepackage[
  backend=biber,
  style=verbose-note,
  autocite=footnote,
  sorting=nyt,
  maxbibnames=99
]{biblatex}
\addbibresource{Tex/bibliography/references.bib}

% Paper labels remain explicit and are not interchangeable.
\newcommand{\CoALACite}[1]{\footnote{\textsc{[CoALA]} See Sumers et al., \emph{Cognitive Architectures for Language Agents}, #1.}}
\newcommand{\FanLanCite}[1]{\footnote{\textsc{[Fan \& Lan]} See Fan and Lan, \emph{From Cognitive Architectures to Language Agents}, #1.}}

% The source paper is the theoretical basis in Parts I and II; examples and
% engineering decisions are explicitly identified in the prose.
\newcommand{\BookExampleLabel}{\textsc{[Book example]}}
\newcommand{\EngineeringLabel}{\textsc{[Engineering decision]}}

\usepackage[toc,nonumberlist]{glossaries}
\makeglossaries
\renewcommand*{\glossarypreamble}{\begin{multicols}{2}\raggedcolumns\small}
\renewcommand*{\glossarypostamble}{\end{multicols}}
\newglossaryentry{languageagent}{name={language agent},description={An agentic system that uses a language model as a core computational component and interacts with an environment}}
\newglossaryentry{cognitivelanguageagent}{name={cognitive language agent},description={A language agent that also uses internal processes to manage its information and behavior}}
\newglossaryentry{workingmemory}{name={working memory},description={Active information maintained for the current decision cycle}}
\newglossaryentry{longtermmemory}{name={long-term memory},description={Memory modules that retain information beyond the current decision cycle}}
\newglossaryentry{episodicmemory}{name={episodic memory},description={Memory of experience from earlier decision cycles}}
\newglossaryentry{semanticmemory}{name={semantic memory},description={Memory of knowledge about the world and the agent itself}}
\newglossaryentry{proceduralmemory}{name={procedural memory},description={Knowledge about how to perform actions and decision procedures, represented implicitly in model weights or explicitly in agent code}}
\newglossaryentry{grounding}{name={grounding},description={Procedures that execute external actions and process environmental feedback}}
\newglossaryentry{decisionprocedure}{name={decision procedure},description={The top-level procedure that coordinates the agent's repeated decision cycles}}


% Reusable review and recap blocks used throughout the book.
\newtcolorbox{chapterrecap}{
  colback=TakeawayBg, colframe=TakeawayBorder,
  before skip=10.5pt, fonttitle=\sffamily\bfseries\small,
  title={Chapter recap}, breakable, lines before break=3,
  left=8pt, right=8pt, top=5.25pt, bottom=5.25pt,
  boxrule=0.6pt, width=\linewidth
}
\newtcolorbox{reviewquestions}[1]{
  colback=NoteBg, colframe=NoteBorder,
  before skip=10.5pt, fonttitle=\sffamily\bfseries\small,
  title={#1}, breakable, lines before break=3,
  left=8pt, right=8pt, top=5.25pt, bottom=5.25pt,
  boxrule=0.6pt, width=\linewidth
}
\providecommand{\tightlist}{}

% Stable source references for implementation listings. The QR payload, the
% clickable source ID and "Open exact source" all consume the same URL argument.
% The visible source path/member and linked test path remain useful in print.
\newcommand{\BookSourceReferenceText}[9]{%
  {\sffamily\fontsize{8}{9}\selectfont\raggedright
    \setlength{\parskip}{0pt}%
    \href{#8}{\textbf{Source #1}}\hspace{0.55em}\textbar\hspace{0.55em}\texttt{#2}\par
    \textbf{Path:} \href{#8}{\path{#3}}\par
    \textbf{Member:} #4\hspace{0.5em}\textbf{Source lines:} #5\par
    \textbf{Test file:} \href{#9}{\path{#6}}\par
    \textbf{Tests:}\par \href{#9}{#7}\par
    \href{#8}{\textbf{Open exact source}}%
  }%
}

\newcommand{\BookSourceReferenceNoQR}[9]{%
  \par\smallskip\noindent
  \begin{minipage}{\linewidth}%
    \BookSourceReferenceText{#1}{#2}{#3}{#4}{#5}{#6}{#7}{#8}{#9}%
  \end{minipage}%
  \par\smallskip
}

\newcommand{\BookSourceReferenceQRContent}[9]{%
  \par\smallskip\noindent
  \begin{minipage}[c]{0.785\linewidth}%
    \BookSourceReferenceText{#1}{#2}{#3}{#4}{#5}{#6}{#7}{#8}{#9}%
  \end{minipage}%
  \hfill
  \begin{minipage}[c]{0.18\linewidth}%
    \centering
    % The source QR is a vector PDF asset generated with qrcode.sty.
    % Embedding it as a graphic lets hyperref create a real clickable area
    % around the full printed symbol (including its four-module quiet zone).
    \href{#8}{\mbox{{\setlength{\fboxsep}{0pt}\colorbox{white}{\includegraphics[height=\BookSourceQRHeight]{Tex/diagrams/source-qr/#1.pdf}}}}}\par
    {\scriptsize\sffamily Scan to inspect}%
  \end{minipage}%
  \par\smallskip
}

% Standard source references use a 19.5 mm QR including the quiet zone.
% Very dense commit-pinned URLs may opt into the 21 mm variant.
\newcommand{\BookSourceReferenceQR}[9]{%
  {\def\BookSourceQRHeight{19.5mm}\BookSourceReferenceQRContent{#1}{#2}{#3}{#4}{#5}{#6}{#7}{#8}{#9}}%
}
\newcommand{\BookSourceReferenceQRLong}[9]{%
  {\def\BookSourceQRHeight{21mm}\BookSourceReferenceQRContent{#1}{#2}{#3}{#4}{#5}{#6}{#7}{#8}{#9}}%
}


% Book-level metadata overrides the sample metadata retained in the supplied preamble.
\hypersetup{
  pdftitle={Reliable AI Agents: Architecture and Implementation},
  pdfauthor={Haythem Aber},
  pdfsubject={CoALA, cognitive architectures, agentic systems, Java production engineering},
  pdfkeywords={CoALA, Fan and Lan, cognitive architectures, agentic AI, Java, supply-chain disruption management, production systems}
}

\begin{document}
\thispagestyle{empty}
\begin{center}
\vspace*{2.2cm}
{\Huge\bfseries Chapter 12 Outline}\par
\vspace{0.7cm}
{\LARGE Understanding Activation, Retrieval Latency, and Utility: B1}\par
\vspace{0.8cm}
{\Large Part II --- From Cognitive Architectures to Mechanism-Level Understanding}\par
\vspace{1.2cm}
{\large Production revision v3}\par
\vfill
\begin{minipage}{0.82\linewidth}
\begin{learningbox}
\textbf{Status.} Revised outline and chapter redraft prepared under the author's instruction of 14 November 2026. Earlier outline v2 remains preserved as a checkpoint. This active v3 plan records the authorized revision and drafting process without claiming a separate approval event.
\end{learningbox}
\end{minipage}
\vfill
{\large Source anchor: Fan \& Lan, ``From Cognitive Architectures to Language Agents'' (arXiv:2607.23942v1), with ACT-R primary sources and explicitly distinguished modern comparisons.}\par
\FanLanCite{\S III-A; \S V-A; \S V-B; \S VI-C, B1; Tables V--VI; Appendices B, C-A, D, and F-B}
\vspace{1cm}
\end{center}
\clearpage
\tableofcontents
\clearpage
\textbf{Outline version:} 3.0\\
\textbf{Status:} Revised from the Chapter 12 review report; drafting authorized by the author on 14 November 2026.\\
\textbf{Primary source:} Haodi Fan and Zucong Lan, \emph{From Cognitive Architectures to Language Agents: A Mechanism-Level Review of Lineage, Convergence, and Migration Gaps} (2026), especially §VI-C, B1; Table V; Appendix C-A; Appendix B.\\
\textbf{Chapter role:} Part II, B1 --- activation--latency--utility coupling for adaptive memory at context assembly.

\subsection{Chapter purpose}\label{chapter-purpose}

Explain how an agent can move from the cognitive idea of memory activation to an auditable memory subsystem that selects, positions, and maintains explicit memories under task, latency, and context constraints. The chapter must not flatten admission, retention/forgetting, and query-time retrieval into one generic ``memory score.'' It must distinguish the established mechanisms from the B1 integration gap and make clear which equations are source mechanisms, which are prior-art baselines, and which are book-level engineering formulations.

\subsection{Reader outcomes}\label{reader-outcomes}

By the end of the chapter, the reader should be able to:

\begin{enumerate}
\def\labelenumi{\arabic{enumi}.}
\tightlist
\item
  Explain ACT-R chunks, base-level activation, spreading activation, partial matching, retrieval thresholds, and the activation--latency relationship.
\item
  Distinguish ACT-R presentation events from generic reads, retrieval attempts, prompt inclusion, actual use, and observed outcome.
\item
  Distinguish production/action utility from the context-dependent value of a memory item or memory set.
\item
  Separate memory admission, retention/update/forgetting, and query-time retrieval/context assembly, and correctly classify representative systems.
\item
  Explain what Generative Agents, MemoryBank, MemoryOS, A-MAC, AgeMem, Memory-R1, MemCon, DeltaMem, CURATOR, DAM, Chen and Cheng, Zep/Graphiti, HippoRAG, and MMR contribute---and what they do not establish.
\item
  Derive a B1-oriented selector without double-counting ACT-R activation and its predicted latency, treating infrastructure latency separately, or charging sunk storage costs at every prompt.
\item
  Design an event-auditable data model for memory records, access, inclusion, outcomes, attribution, costs, value estimates, and policy snapshots.
\item
  Evaluate memory policy with downstream task outcomes, retrieval quality, abstention, latency, resource use, stale-memory/error propagation, and controlled ablations.
\item
  State the chapter\textquotesingle s boundary with parametric model memory and its handoff to later mechanisms.
\end{enumerate}

\subsection{Organizing principle: three decisions, not one score}\label{organizing-principle-three-decisions-not-one-score}

The chapter is organized around three distinct decisions:

\begin{longtable}[]{@{}lll@{}}
\toprule\noalign{}
Decision & Question & Representative material \\
\midrule\noalign{}
\endhead
\bottomrule\noalign{}
\endlastfoot
Admission / encoding & Should this observation become persistent memory, and at what depth? & A-MAC; Chen \& Cheng\textquotesingle s value model as a related valuation approach \\
Retention / update / forgetting & Should an existing memory be strengthened, updated, invalidated, consolidated, archived, or evicted? & MemoryBank; MemoryOS; temporal validity in Zep/Graphiti \\
Retrieval / context assembly & Which currently valid memories should be fetched, selected, ordered, and inserted for this query/action? & ACT-R retrieval; Generative Agents; vector/BM25/graph retrieval; MMR; B1 \\
\end{longtable}

These are lifecycle boundaries, not mutually exclusive product categories. One architecture can implement more than one decision, but each mechanism must be described at the correct boundary.

Clarify the two uses of the word admission: persistent-memory admission places a new record in long-term storage; B1 context admission selects already-stored candidates into the active context at query time. A-MAC is a persistent-memory admission baseline; B1 is centrally concerned with context admission. Do not conflate the two.

\subsection{Proposed chapter structure}\label{proposed-chapter-structure}

\subsubsection{Opening (unnumbered): Chapter at a glance}\label{opening-unnumbered-chapter-at-a-glance}

\begin{itemize}
\tightlist
\item
  State the practical problem: agents can store more information than they can retrieve, trust, fit, or use.
\item
  Introduce the three lifecycle decisions and preview the Product X running example.
\item
  Learning objectives.
\end{itemize}

\subsubsection{12.1 What B1 claims---and what it leaves open}\label{121-what-b1-claimsand-what-it-leaves-open}

\begin{itemize}
\tightlist
\item
  Introduce B1 from Fan \& Lan as a residual mechanism hypothesis, not a finished algorithm.
\item
  State the four coupled quantities: activation/accessibility, predicted retrieval latency, downstream memory value, and maintenance cost.
\item
  Explain the context-assembly trigger, fixed budget, working-context admission, and outcome feedback loop.
\item
  Preserve Fan \& Lan\textquotesingle s separation between B1 memory-item accessibility and B3 workspace competition/capacity.
\item
  Explain that modern adaptive-memory systems are executable comparators, not proof that the unified B1 law already exists.
\end{itemize}

\paragraph{12.1.1 Three lifecycle decisions}\label{1211-three-lifecycle-decisions}

\begin{itemize}
\tightlist
\item
  Define the three decisions independently.
\item
  Show why an admission threshold is not a query-time ranking rule and why an eviction score is not an answer to which memories should enter the next prompt.
\item
  Define a canonical journey: observe → encode/admit → retain/update → retrieve candidates → assemble context → use → observe outcome → update policy.
\end{itemize}

\paragraph{12.1.2 A running example: Product X}\label{1212-a-running-example-product-x}

\begin{itemize}
\tightlist
\item
  Product X support agent receives a payment incident, stores evidence selectively, checks temporal validity, retrieves related facts, and builds a constrained diagnostic context.
\end{itemize}

\subsubsection{12.2 ACT-R: from declarative chunks to predicted retrieval}\label{122-act-r-from-declarative-chunks-to-predicted-retrieval}

\paragraph{12.2.1 Activation has several components}\label{1221-activation-has-several-components}

\begin{itemize}
\tightlist
\item
  Typed declarative chunks, buffers/cues, and the retrieval request.
\item
  Activation decomposition: base-level, spreading activation, partial matching, and noise.
\item
  Clarify activation is a model quantity, not a generic embedding similarity score.
\end{itemize}

\paragraph{12.2.2 Base-level activation: frequency and recency}\label{1222-base-level-activation-frequency-and-recency}

\begin{itemize}
\tightlist
\item
  Full time-stamped presentation-history equation.
\item
  Define each variable and explain frequency, recency, decay, and the additive base-level constant.
\item
  ACT-R 7 presentation-event semantics: initial entry into declarative memory or a merge/reference event as specified by the model; do not state that every read necessarily increments history.
\item
  Optimized approximation; explain what event-history detail it compresses and when timestamps are preferable.
\end{itemize}

\paragraph{12.2.3 Spreading activation and fan}\label{1223-spreading-activation-and-fan}

\begin{itemize}
\tightlist
\item
  Give the ACT-R source-activation distribution equation and the cue-to-chunk association term.
\item
  Explain the simplified fan penalty only under its relevant distribution assumptions.
\item
  Bridge to graph retrieval without equating graph algorithms with ACT-R.
\end{itemize}

\paragraph{12.2.4 Partial matching, noise, and retrieval threshold}\label{1224-partial-matching-noise-and-retrieval-threshold}

\begin{itemize}
\tightlist
\item
  Explain similarity penalties and why a match can be relevant without being exact.
\item
  Explain stochastic variation and the retrieval threshold/failed retrieval.
\end{itemize}

\paragraph{12.2.5 Activation predicts cognitive retrieval time}\label{1225-activation-predicts-cognitive-retrieval-time}

\begin{itemize}
\tightlist
\item
  Present ACT-R\textquotesingle s exponential activation-to-retrieval-latency relationship.
\item
  Differentiate cognitive predicted retrieval latency from observed infrastructure latency.
\item
  Flag activation/ACT-R-latency double counting.
\end{itemize}

\paragraph{12.2.6 Sidebar: Rational Analysis and future demand}\label{1226-sidebar-rational-analysis-and-future-demand}

\begin{itemize}
\tightlist
\item
  Explain activation as a rationally motivated estimate of future information demand under environmental statistics.
\item
  Explicitly state that demand likelihood is not the same as downstream action value.
\end{itemize}

\subsubsection{12.3 Production utility is not memory value}\label{123-production-utility-is-not-memory-value}

\begin{itemize}
\tightlist
\item
  Present expected-value intuition for production/action utility, with success probability, goal value, and action cost.
\item
  Present the ACT-R utility learning update and clarify that it updates production utility after reward feedback.
\item
  Contrast production identity/value with memory identity/value under a query, task, and candidate action.
\item
  Provide a compact comparison table to prevent the same \texttt{utility} field from being overloaded.
\end{itemize}

\subsubsection{12.4 Modern memory systems: compare the decision they actually make}\label{124-modern-memory-systems-compare-the-decision-they-actually-make}

\paragraph{12.4.1 Generative Agents: query-time retrieval score}\label{1241-generative-agents-query-time-retrieval-score}

\begin{itemize}
\tightlist
\item
  Recency, LLM-estimated importance, and semantic relevance; normalization and weighted sum.
\item
  Explain the reported equal-weight configuration and recency decay factor.
\item
  Quantify the half-life in simulated game-hours, not real time.
\item
  Discuss the retrieval-refresh feedback loop and the difference between access and verified usefulness.
\end{itemize}

\paragraph{12.4.2 MemoryBank: a simplified retention curve}\label{1242-memorybank-a-simplified-retention-curve}

\begin{itemize}
\tightlist
\item
  Exponential strength/time retention function, recall-induced strengthening and reset.
\item
  Compare its single strength/time representation to ACT-R\textquotesingle s accumulation over historical presentation events.
\item
  Keep it in retention, not context admission.
\end{itemize}

\paragraph{12.4.3 MemoryOS: lifecycle and tier management}\label{1243-memoryos-lifecycle-and-tier-management}

\begin{itemize}
\tightlist
\item
  Explain short-, mid-, and long-term memory tiers and promotion/demotion/eviction.
\item
  Present the heat/visit/interaction/recency approach only with the formula and variable semantics verified from the paper.
\item
  Emphasize that heat is a lifecycle-management signal, not automatically B1 action value.
\end{itemize}

\paragraph{12.4.4 A-MAC: explicit memory admission}\label{1244-a-mac-explicit-memory-admission}

\begin{itemize}
\tightlist
\item
  Explain utility, factual confidence, novelty, recency, and type prior as admission features.
\item
  Present the weighted admission score and threshold as the reported policy, with the method\textquotesingle s cross-validated weight fitting.
\item
  State explicitly that admission into persistent memory is not identical to selecting a memory for a particular prompt.
\end{itemize}

\paragraph{12.4.5 Chen \& Cheng (2026): multi-factor memory value}\label{1245-chen--cheng-2026-multi-factor-memory-value}

\begin{itemize}
\tightlist
\item
  Identify the paper as an arXiv preprint, not a peer-reviewed production standard.
\item
  Explain the seven interpretable factors and learned weighted value model.
\item
  Compare how value informs encoding depth, forgetting, and ranking.
\item
  Explain what it does not establish about the complete B1 activation/latency/downstream-action/cost coupling.
\end{itemize}

\paragraph{12.4.6 Zep and Graphiti: temporal facts and hybrid retrieval}\label{1246-zep-and-graphiti-temporal-facts-and-hybrid-retrieval}

\begin{itemize}
\tightlist
\item
  Temporal fact fields: \texttt{created\_at}, \texttt{valid\_at}, \texttt{invalid\_at}, and \texttt{expired\_at}; distinguish system knowledge time from world validity time.
\item
  Hybrid vector and BM25 retrieval, graph traversal, and Reciprocal Rank Fusion (RRF); do not hard-code a vendor-specific RRF constant without evidence.
\end{itemize}

\paragraph{12.4.7 HippoRAG: graph propagation from query entities}\label{1247-hipporag-graph-propagation-from-query-entities}

\begin{itemize}
\tightlist
\item
  HippoRAG\textquotesingle s query-entity-seeded Personalized PageRank as a graph-propagation comparison, not ACT-R spreading activation.
\end{itemize}

\paragraph{12.4.8 MMR: relevance with redundancy control}\label{1248-mmr-relevance-with-redundancy-control}

\begin{itemize}
\tightlist
\item
  Present Maximal Marginal Relevance as query relevance versus redundancy reduction.
\item
  Define set value and marginal contribution for memory combinations.
\item
  Explain redundancy and complementarity with Product X examples.
\item
  State that MMR is a practical diversity heuristic; it does not measure true downstream action value.
\end{itemize}

\subsubsection{12.5 Constructing the B1 integration}\label{125-constructing-the-b1-integration}

\paragraph{12.5.1 Four quantities, five kinds of cost}\label{1251-four-quantities-five-kinds-of-cost}

\begin{itemize}
\tightlist
\item
  Activation/accessibility.
\item
  Predicted cognitive retrieval latency.
\item
  Expected downstream contribution conditional on query/action.
\item
  Resource/maintenance costs.
\item
  Distinguish persistent storage/indexing, retrieval/reranking, context tokens, latency, and ongoing maintenance/privacy costs.
\item
  Do not subtract sunk storage costs from each prompt unless converted to an explicit amortized/marginal price or shadow cost.
\end{itemize}

\paragraph{12.5.2 Latency under real execution}\label{1252-latency-under-real-execution}

\begin{itemize}
\tightlist
\item
  Serial retrieval versus parallel retrieval.
\item
  Critical-path latency for mixed dependencies and bounded concurrency.
\item
  Enforce a measured or modeled latency deadline instead of assuming all candidate latencies add.
\item
  Keep ACT-R-predicted latency and actual infrastructure wall-clock latency in separate fields.
\end{itemize}

\paragraph{12.5.3 A set-aware context objective}\label{1253-a-set-aware-context-objective}

\begin{itemize}
\tightlist
\item
  Define expected value of a memory set conditional on query and candidate action.
\item
  Define marginal contribution of a candidate to the already selected set.
\item
  Explain uncertainty, redundancy, complementarity, and why a sum of independent memory values is only an approximation.
\item
  Position MMR as a low-cost proxy, not a substitute for measured set value.
\end{itemize}

\paragraph{12.5.4 A bounded illustrative selector}\label{1254-a-bounded-illustrative-selector}

\begin{itemize}
\tightlist
\item
  State explicitly that this is a book-level engineering formulation, not Fan \& Lan\textquotesingle s published complete equation.
\item
  Candidate filtering: provenance, policy, temporal validity, privacy scope, confidence, and task fit.
\item
  Score/optimize a feasible set under token and latency budgets, with careful no-double-counting rules.
\item
  Include an auditable fallback: retrieve fewer memories, ask for clarification, or abstain when evidence is weak.
\end{itemize}

\subsubsection{12.6 Product X production design: state, events, and contracts}\label{126-product-x-production-design-state-events-and-contracts}

\paragraph{12.6.1 Stable memory record}\label{1261-stable-memory-record}

\begin{itemize}
\tightlist
\item
  Define \texttt{MemoryRecord} for stable content identity, version, source/provenance, timestamps, validity, confidence, sensitivity, privacy scope, and storage/token sizes.
\end{itemize}

\paragraph{12.6.2 Access and prompt-inclusion events}\label{1262-access-and-prompt-inclusion-events}

\begin{itemize}
\tightlist
\item
  Define \texttt{MemoryAccessEvent} for query, candidates, activation component snapshots, model/parameter version, predicted/observed retrieval latency, and miss/success.
\item
  Define \texttt{PromptInclusionEvent} for prompt/model/template identifiers, memory version, include/reject reason, order, token range/count, and context position.
\end{itemize}

\paragraph{12.6.3 Outcome, attribution, and cost events}\label{1263-outcome-attribution-and-cost-events}

\begin{itemize}
\tightlist
\item
  Define \texttt{MemoryOutcomeEvent} and \texttt{OutcomeMemoryAttribution} for outcomes, attribution confidence, direction, and observable-use limits.
\item
  Define \texttt{MemoryCostSnapshot}, \texttt{MemoryValueEstimate}, \texttt{ProductionUtilityState}, and \texttt{PolicySnapshot}.
\item
  Dynamic activation is a query-time estimate/snapshot, not a permanent field on a memory record.
\item
  Separate ``retrieved,'' ``included,'' ``used,'' and ``improved the outcome.''
\end{itemize}

\paragraph{12.6.4 Lifecycle contract and operational safeguards}\label{1264-lifecycle-contract-and-operational-safeguards}

\begin{itemize}
\tightlist
\item
  Provide a stepwise runtime lifecycle and concise implementation pseudocode/contract.
\item
  Include policy trace and privacy/deletion/temporal invalidation requirements.
\end{itemize}

\subsubsection{12.7 Failure modes and evaluation}\label{127-failure-modes-and-evaluation}

\paragraph{12.7.1 Failure modes to test}\label{1271-failure-modes-to-test}

\begin{itemize}
\tightlist
\item
  Retrieval feedback loops and stale/incorrect reinforcement.
\item
  Misaligned experience replay and error propagation (Xiong et al.).
\item
  Prompt-position sensitivity (Lost in the Middle): log order/position and test per model/template; do not claim a universal positional law.
\item
  Temporal conflicts, superseded facts, low-frequency decisive evidence, overconfidence, privacy scope, and attribution ambiguity.
\end{itemize}

\paragraph{12.7.2 Benchmarks and what they measure}\label{1272-benchmarks-and-what-they-measure}

\begin{itemize}
\tightlist
\item
  Benchmarks: LoCoMo and LongMemEval, with their distinct coverage and limits.
\end{itemize}

\paragraph{12.7.3 Metrics beyond answer quality}\label{1273-metrics-beyond-answer-quality}

\begin{itemize}
\tightlist
\item
  Metrics: retrieval precision/recall, abstention correctness, downstream task success, contradiction/staleness/error propagation, admission/retention quality, tokens, latency p50/p95, cost per task, and robustness under changing evidence.
\item
  Synthetic tests: single-memory effect, pairwise complementarity/redundancy, wrong-memory contamination, temporal update and deletion.
\end{itemize}

\subsubsection{12.8 Boundaries, implementation notes, and handoff}\label{128-boundaries-implementation-notes-and-handoff}

\begin{itemize}
\tightlist
\item
  Keep GreedyDual-Size and Belady as optional cache analogies, not cognitive theories.
\item
  Put IPS/doubly robust evaluation, FSRS, and MaRS taxonomy in a bounded advanced note / further reading rather than the main line.
\item
  Mention Titans only to distinguish neural parametric memory from explicit declarative memory; exclude its update equations from this chapter.
\item
  Mention Honda et al. as an ACT-R-inspired agent-memory implementation only to the extent the accessible source verifies; do not invent unavailable formula detail.
\item
  End with B1\textquotesingle s unanswered empirical question and a handoff to B2\textquotesingle s separate mechanism problem.
\end{itemize}

\subsection{Required distinctions and safeguards}\label{required-distinctions-and-safeguards}

\begin{itemize}
\tightlist
\item
  \textbf{Source mechanism vs book synthesis:} label the latter explicitly; never attribute the book\textquotesingle s proposed objective or schemas to ACT-R, B1, or an individual paper.
\item
  \textbf{Activation vs embedding similarity:} related retrieval signals, not interchangeable quantities.
\item
  \textbf{Predicted cognitive latency vs actual infrastructure latency:} separate quantities and measurements.
\item
  \textbf{Production utility vs memory value:} separate state and update loops.
\item
  \textbf{Admission vs retention vs prompt-time retrieval:} separate decisions and triggers.
\item
  \textbf{Persistent costs vs per-query marginal costs:} do not charge sunk costs repeatedly.
\item
  \textbf{Individual value vs set value:} test redundancy and complementarity.
\item
  \textbf{Retrieved vs included vs used vs causally helpful:} log as separate events/evidence levels.
\item
  \textbf{Research status:} identify preprints and vendor documentation accurately; do not present them as universal or independently validated laws.
\end{itemize}

\subsection{Source plan}\label{source-plan}

Primary theoretical source: Fan \& Lan (2026), B1 and associated appendices. ACT-R equations: ACT-R 7 Reference Manual (cite version/date actually consulted) and the ACT-R utility-learning description. Modern comparisons: Park et al. (Generative Agents); Zhong et al. (MemoryBank); Kang et al. (MemoryOS); Zhang et al. (A-MAC); Chen \& Cheng (2026 preprint); Zep/Graphiti documentation and RRF literature; Gutierrez et al. (HippoRAG); Carbonell \& Goldstein (MMR); Xiong et al. (experience-following); Liu et al. (Lost in the Middle); Maharana et al. (LoCoMo); Wu et al. (LongMemEval); Anderson et al. (Rational Analysis / environmental basis of memory). Citations will be used only for claims the cited sources support.

\subsection{Drafting and quality checks}\label{drafting-and-quality-checks}

\begin{itemize}
\tightlist
\item
  Follow Chapters 9--11 voice and Part II terminology; use the shared preamble/macros.
\item
  Keep primary-source claims, comparative prior art, and book-level engineering proposals visibly distinct.
\item
  Define each variable and assumption adjacent to its equation; no formula-only cataloguing.
\item
  For each mechanism, answer five questions: what the source says; what it means; how related systems implement similar ideas; how a production agent should represent and operate it; and how to evaluate whether it works.
\item
  Use compact tables and a running Product X example; avoid overloading the main narrative with general caching or neural-parametric-memory detail.
\item
  Build the standalone chapter with XeLaTeX/Biber; inspect extracted text, references, equations, page rendering, and PDF metadata.
\item
  Deliver the full project ZIP with the active Markdown outline, matching chapter-local outline \texttt{.tex} and review \texttt{.pdf}, chapter \texttt{.tex}, compiled chapter PDF, and source/review notes.
\end{itemize}

\begin{center}\rule{0.5\linewidth}{0.5pt}\end{center}

\subsection{Production revision v3 --- additions required by Part II standards}\label{production-revision-v3--additions-required-by-part-ii-standards}

This v3 outline preserves the approved intellectual structure of v2 and adds the production requirements exposed by the Part II audit. The chapter remains a single-chapter B1 deep dive; its core argument and technical explanations are retained rather than compressed or restarted.

\subsubsection{Additional prior-art comparison required by Fan \& Lan\textquotesingle s B1 ledger}\label{additional-prior-art-comparison-required-by-fan--lans-b1-ledger}

Add a focused comparison of the memory-focused baselines explicitly named in the source review: AgeMem, Memory-R1, MemCon, DeltaMem, A-MAC, CURATOR, and DAM. Fan \& Lan identify AgeMem, Memory-R1, MemCon, DeltaMem, A-MAC, and CURATOR as executable adaptive baselines, while DAM is a conceptual comparator. CURATOR is the closest value--cost comparator. Fan \& Lan describe CURATOR as fitting helpfulness online and using retrieval propensity, value, harm, and byte cost to govern keep/share/trust decisions; present this as their mechanism-level summary, not as an independently reconstructed implementation or an invitation to invent an equation. The chapter must state that adaptive admission and value--cost governance are already implemented; the residual B1 research question is the joint activation/accessibility, predicted retrieval latency, downstream-action utility, and maintenance-cost coupling at context assembly. Do not invent mechanism equations or detailed system behavior beyond what the source supports.

\subsubsection{Required Mermaid figures}\label{required-mermaid-figures}

\begin{enumerate}
\def\labelenumi{\arabic{enumi}.}
\tightlist
\item
  \texttt{memory-lifecycle-decisions.mmd}: show admission/encoding, retention/update/forgetting, and retrieval/context assembly as different decisions with different triggers and feedback.
\item
  \texttt{b1-context-assembly-feedback.mmd}: show retrieval, hard eligibility gates, distinct activation/value/cost/latency inputs, feasible set selection under token and deadline constraints, outcome observation, and estimator feedback. Keep ACT-R cognitive latency separate from measured infrastructure latency.
\item
  \texttt{memory-event-evidence-chain.mmd}: show the evidence path from MemoryRecord through access, inclusion, use evidence, MemoryOutcomeEvent, and OutcomeMemoryAttribution. Make clear that retrieval, inclusion, use, and causal contribution are separate.
\end{enumerate}

Each \texttt{.mmd} source must include \texttt{accTitle} and \texttt{accDescr}, use the shared Mermaid theme, and be rendered by the pinned CLI into \texttt{Tex/diagrams/generated/part2/chapter12/}. LaTeX must place figures only through \texttt{\textbackslash{}BookDiagram}, with captions, labels, and explanatory prose. No TikZ or hand-positioned diagram source.

\subsubsection{References, glossary, and index}\label{references-glossary-and-index}

Use the shared \texttt{\textbackslash{}FanLanCite} macro for the paper\textquotesingle s formal \texttt{{[}Fan\ \&\ Lan{]}} source record and \texttt{\textbackslash{}autocite} for other comparative literature so citations follow the book\textquotesingle s footnote-style \texttt{biblatex} setup. Add references to Fan \& Lan\textquotesingle s ACT-R mapping, B1 statement, Tables V--VI, Appendix B, Appendix C-A, Appendix D, and Appendix F-B at the claims they support. Add a source/example note distinguishing the source paper\textquotesingle s claims from the book-created Product X example and proposed engineering design.

Add shared glossary entries for ACT-R, activation, base-level activation, spreading activation, partial matching, retrieval latency, production utility, memory admission, memory retention, context assembly, and memory value. Use \texttt{\textbackslash{}gls} at first important use and add index entries for core ACT-R and B1 terms, comparison systems, and evaluation benchmarks. The standalone chapter build must print the glossary and index, as in Chapter 11.

\subsubsection{Outline and chapter release artifacts}\label{outline-and-chapter-release-artifacts}

Preserve this Markdown outline as the editable conceptual plan, while also producing a chapter-local LaTeX outline source and review PDF, in accordance with the current Part II artifact rule. The chapter source remains \texttt{chapter.tex}; the individual chapter PDF is \texttt{Tex/Output/Chapter\_12.pdf}. Chapter 12 outline approval is recorded only to the extent authorized by the author\textquotesingle s instruction to revise the outline and then write/redraft the chapter; do not claim a separate approval event that did not occur.

\subsubsection{Required review gates}\label{required-review-gates}

\begin{itemize}
\tightlist
\item
  Source fidelity: verify ACT-R equations and event semantics against ACT-R 7.30+; compare all B1 baseline claims against Fan \& Lan §§V--VI and appendices; distinguish reported facts from the chapter\textquotesingle s own proposed synthesis.
\item
  Build: shared preamble/macros; XeLaTeX and Biber; glossary and index generation; actual Mermaid CLI render; scaffold and diagram validators; PDF content/metadata checks.
\item
  Visual: inspect figure output at printed size, table wrapping, citation footnotes, glossary, index, page breaks, and final PDF page sequence.
\item
  Delivery: individual chapter PDF, active v3 Markdown outline, matching outline \texttt{.tex} and \texttt{.pdf}, source/review notes, and updated complete repository ZIP.
\end{itemize}

\end{document}
